% §12.5 Bibliographic Notes
%
% Monograph register: the relation manuscript, Codd, the similarity-join
% literature. The "exercise 10" reference is repointed to rem:bloom-equijoin.
% Citations towell-bernoulli-relations, codd-1970, sarawagi-kirpal-2004,
% broder-1997, indyk-motwani-1998. Cross-refs sec:relation-composition,
% thm:join-rates, sec:approximate-joins, rem:bloom-equijoin. DEFINED:
% sec:bib-notes-ch12.

\section*{Bibliographic Notes}
\addcontentsline{toc}{section}{Bibliographic Notes}
\label{sec:bib-notes-ch12}

\paragraph{The approximate-relation framework.}
The relation-level lift of the chapter-5 framework follows
\textcite{towell-bernoulli-relations}, the primary source for the approximate
relation as a $\Bernoulli{\BoolType}$ on the pair universe and for the join,
projection, and selection rate rules of \Cref{sec:relation-composition}. The
chapter reads those rules as the chapter-5 and chapter-6 frameworks specialized
to pairs, making \Cref{thm:join-rates} a corollary of the chapter-6 intersection
theorem with cross-relation independence in the role of cross-structure
independence.

\paragraph{The classical relational model.}
The model of a relation as a subset of a Cartesian product, and the operators the
chapter approximates, originate with \textcite{codd-1970}, the foundation of the
relational-database tradition. The chapter borrows the operators wholesale; the
only additions are the approximate lift and its rate analysis.

\paragraph{Probabilistic and uncertain databases.}
A separate tradition also attaches uncertainty to relational data.
\textcite{suciu-2011} survey probabilistic databases, which assign each tuple a
probability of belonging to the relation and answer queries over the induced
distribution of possible worlds; the systems Trio \parencite{widom-2005-trio} and
MayBMS \parencite{huang-2009-maybms} implement the model. The approximate relation
here is a different object. Its uncertainty is not a probability that a stored
fact is true but a per-query error rate of a fixed approximate structure: the
relation is exact and known, and the randomness lives in the membership test,
with the rates $\fprate$ and $\fnrate$ describing how a Bernoulli trial reports a
definite truth. Possible-worlds semantics ask which tuples are probably present;
the approximate-relation semantics ask, for a given tuple, how reliably its
membership comes back. The two are complementary rather than competing, one
modeling uncertain data tested exactly, the other certain data tested
approximately.

\paragraph{Similarity-join algorithms.}
The algorithmic instance of \Cref{sec:approximate-joins} combines
\textcite{broder-1997}'s MinHash with \textcite{indyk-motwani-1998}'s
locality-sensitive hashing, systematized as a database operator for
set-similarity joins by \textcite{sarawagi-kirpal-2004}; the many later variants
(prefix, positional, and length filtering) are each an alternative
implementation of the same approximate relation $\tilde S_\theta$, with rates by
the same chapter-12 machinery. \Cref{thm:join-rates} commits to no particular
sketch: the chapter-10 MinHash bound is used only because the book already
developed it, and the Bloom equijoin of \Cref{rem:bloom-equijoin} is another
instance with different per-pair rates.
